\documentclass[10pt,a4paper]{amsart}
\usepackage[margin=19mm]{geometry}
\usepackage{amsmath,amssymb,mathtools}
\usepackage[scr=rsfs,cal=euler]{mathalfa}
\usepackage{xfrac}
\usepackage[unicode,hidelinks,bookmarks=false]{hyperref}
\newcommand{\ie}{i.e.\ }
\newcommand{\cf}{cf.\ }
\newcommand{\resp}{respectively\ }
\newcommand{\Subsection}{\subsection}
\providecommand{\nobreakdash}{\nobreak\hbox{-}}
\setlength{\emergencystretch}{2em}
\multlinegap0pt
\usepackage{tikz-cd}
\usetikzlibrary{decorations.pathreplacing}
\usepackage{physics}
\usepackage{multicol,multirow,array}
\usepackage{cleveref}
\newcommand{\eventorus}{%
    \tikz[baseline=-0.5ex, scale=0.35]{
        \draw[fill=PineGreen, fill opacity=0.3, line width=1pt]
            (-0.5,0) -- (-0.5,4) -- (3.5,4) -- (3.5,0) -- cycle;
    }%
}
\newcommand{\oddklein}{%
    \tikz[baseline=-0.5ex, scale=0.35]{
        \draw[fill=PineGreen, fill opacity=0.3, line width=1pt]
            (6,0) -- (6,4) -- (10,4) -- (14,0) -- cycle;
    }%
}
\newtheorem{theorem}{Theorem}[section]
\newtheorem{theoremletter}{Theorem}
\renewcommand{\thetheoremletter}{\Alph{theoremletter}}
\newtheorem{proposition}[theorem]{Proposition}
\newtheorem{corollary}[theorem]{Corollary}
\newtheorem{lemma}[theorem]{Lemma}
\newtheorem{definition}[theorem]{Definition}
\newtheorem*{recall}{Recall}
\newtheorem*{remark}{Remark}
\newtheorem*{remarks}{Remarks}
\newtheorem{example}[theorem]{Example}
\newtheorem{assumption}{Assumption}
\newtheorem*{question}{Question}
\renewcommand\qedsymbol{$\blacksquare$}
\newcommand{\ZZ}{\mathbb{Z}}
\renewcommand{\dd}{\textup{d}}
\renewcommand{\d}{\textup{d}}
\newcommand{\ft}{\mathfrak{t}}
\newcommand{\ftz}{\ft_\ZZ}
\newcommand{\fg}{\mathfrak{g}}
\newcommand{\fs}{\mathfrak{s}}
\newcommand{\red}{\textup{red}}
\newcommand{\lift}{\textup{lift}}
\newcommand{\codim}{\textup{codim}}
\newcommand{\IR}{\mathbb{R}}
\newcommand{\IC}{\mathbb{C}}
\newcommand{\IZ}{\mathbb{Z}}
\newcommand{\FF}{\mathbb{F}}
\newcommand{\IT}{\mathbb{T}}
\newcommand{\ICP}{\mathbb{CP}}
\newcommand{\woz}{\backslash\{0\}}
\newcommand{\pr}{\textup{pr}}
\newcommand{\Hh}{\mathcal{H}}
\newcommand{\std}{\textup{std}}
\newcommand{\dR}{\textup{dR}}
\newcommand{\FS}{\textup{FS}}
\newcommand{\R}{\mathbb{R}}
\newcommand{\C}{\mathbb{C}}
\newcommand{\Z}{\mathbb{Z}}
\newcommand{\N}{\mathbb{N}}
\newcommand{\F}{\mathbb{F}}
\newcommand{\vf}[1]{\mathfrak{X}\left( #1 \right)}
\newcommand{\fgz}{\mathfrak{g}_\mathbb{Z}}
\newcommand{\Oo}{\mathcal{O}}
\newcommand{\Aut}{\textup{Aut}}
\newcommand{\id}{\textup{id}}
\newcommand{\Lie}{\textup{Lie}}
\newcommand{\diff}{\textup{Diff}}
\renewcommand{\sp}{\textup{Symp}}
\newcommand{\resc}{{\textup{res},c}}
\newcommand{\cutc}{{\textup{cut},c}}
\newcommand{\cut}{\textup{cut}}
\newcommand{\im}{\textup{Im}}
\newcommand{\hor}{\textup{hor}}
\newcommand{\tot}{\textup{tot}}
\newcommand{\Ham}{\textup{Ham}}
\newcommand{\iso}{\textup{iso}}
\newcommand{\joe}[1]{\textcolor{blue}{\footnotesize#1}}
\newcommand{\reto}[1]{\textcolor{teal}{\footnotesize#1}}
\renewcommand{\labelenumi}{(\theenumi)}
\begin{document}
\title[The spread construction and non-uniqueness of visible Lagran]{The spread construction and non-uniqueness of visible Lagrangians}
\author{Joé Brendel, Reto Kaufmann}
\date{}
\maketitle
\noindent\emph{Source and licence.} The spread construction and non-uniqueness of visible Lagrangians. Version arXiv:2607.14010v2 (CC BY 4.0). This material is distributed under the Creative Commons Attribution 4.0 International licence, http://creativecommons.org/licenses/by/4.0/, as recorded in the arXiv OAI licence metadata for this exact version and shown on the abstract page https://arxiv.org/abs/2607.14010v2. Full rights notice: licenses/arxiv-2607.14010-CC-BY-4.0.txt.\par\smallskip
\noindent\textit{Editorial scope.} Counted unit: the original \texttt{proposition} environment \texttt{prop:DHmodel} at source lines 674--696 together with its complete proof at lines 697--731; the delivered document keeps source lines 663--732 so that every referenced label and every citation resolves inside it. Native editable source is the paper's own concatenated TeX; the AMS wrapper is editorial formatting only. No publisher class, upstream script or unrelated artwork is redistributed.
\section{The Spread Construction} \label{sec:spreadconstruction} 
The goal of this section is to define and discuss the spread construction. We start by recollecting facts surrounding the Duistermaat--Heckman normal form theorem for symplectic quotients, before moving to the spread construction and then discussing its compatibility with symplectic cuts. The last subsection is dedicated to lifting torus actions from the base to Duistermaat-Heckman model spaces.

\subsection{The Duistermaat--Heckman Theorem}
The name \emph{Duistermaat--Heckman theorem} usually refers to the statement that the cohomology class of the reduced symplectic form varies linearly with the reduction level. For our purpose, we need the geometric refinement describing the symplectic forms themselves. However, this description depends on the choice of a connection on the reduction bundle. Instead of passing to cohomology, we will impose a condition on this connection. 
\begin{theorem}[Duistermaat--Heckman Theorem for $S^1$-actions]
    Let $(M,\omega,H)$ be a symplectic manifold equipped with a Hamiltonian $S^1$-action generated by $H:M\to \IR$. Assume that $S^1$ acts freely on $H^{-1}(0)$ and let $(B,\omega_B)$ denote the symplectic quotient. Then there exist $a<0<b$ such that $M$ admits a symplectic quotient $(M_c,\omega_c)$ at all $c\in (a,b)$ and $$(M_c,\omega_c)\cong (B,\omega_B + c\, F),$$ where $F\in \Omega^2(B)$ is the curvature form of a connection on the principal circle bundle $H^{-1}(0) \to B.$
\end{theorem}

The proof of the Duistermaat--Heckman theorem proceeds by identifying a neighborhood of the regular level set $H^{-1}(0)$ with a model symplectic manifold whose reduced spaces can be described explicitly. We call this model symplectic manifold the \emph{Duistermaat--Heckman model}.
Its symplectic quotients are easy to compute and exhibit the linear variation of the reduced symplectic forms. In the following proposition, we introduce this model and summarize its main properties. For simplicity, we restrict to the case that we are interested in, that is, we assume that the curvature is proportional to the symplectic form, resulting in a linear scaling of the symplectic form for the quotients.

\begin{proposition}[Duistermaat--Heckman Model]
\label{prop:DHmodel}
Let $(B,\omega_B)$ be a symplectic manifold, $\lambda\in \IR$ a real constant and $Z\to B$ a principal $S^1$-bundle with connection form $\alpha\in \Omega^1(Z)$ whose curvature is $F=\lambda\omega_B.$
Let $I\subset\R$ be an open interval and set $$(M,\omega):=\left(Z\times I,(1+\lambda r)\pi^*\omega_B+\dd r\wedge\alpha\right),$$ where $r:M\to I$ is the projection to the second factor and $\pi:M\to B$ is the composition of the projection $M\to Z$ and the bundle map $Z\to B$. Then:
\begin{enumerate}
\item $\omega$ is symplectic on $\{1+\lambda r\neq0\}$;
\item The induced $S^1$-action on $(M,\omega)$ is Hamiltonian with moment map
$r:M\to \IR$;
\item For every regular value $c$ of $r$ (equivalently $1+\lambda c\neq 0$), the reduced space $(M_c,\omega_c)$ defined by
\begin{equation*}
    \begin{tikzcd}
        r^{-1}(c) \arrow[r,hook,"i_c"] \arrow[d,two heads,"\pi_c"',"/S^1"] & (M,\omega) \\
        (M_c,\omega_c)
    \end{tikzcd}
\end{equation*}
and $\pi_c^*\omega_c = i_c^*\omega$ is a symplectic manifold. There exists a
diffeomorphism $\chi_c:M_c \longrightarrow B$ such that
\begin{equation*}
    \chi_c \circ \pi_c = \pi \circ i_c \qq{and} \omega_c = (1+\lambda c) \chi_c^*\omega_B
.
\end{equation*}
\end{enumerate}
\end{proposition}

\begin{proof}
    \begin{enumerate}
        \item Using that $\dd \alpha =\pi^* F$ and $F=\lambda \omega_B$, we can write
        \begin{equation*}
            (1+\lambda r)\pi^*\omega_B + \dd r\wedge \alpha= \pi^*\omega_B + \dd (r \alpha)
        \end{equation*}
        so $\omega$ is closed. To prove nondegeneracy we use the splitting of the tangent bundle defined by the connection form $\alpha$: If $\partial_\theta\in \mathfrak{X}(Z)$ is the infinitesimal generator of the principal $S^1$-action, then 
        \begin{equation*}
            T(Z\times I) = \ker(\alpha) \oplus \expval{\partial_\theta} \oplus \expval{\partial_r}.
        \end{equation*}
        $\pi^*\omega_B$ is a nondegenerate form on the horizontal part $\ker(\alpha)\cong\pi^*(TB)$, while $\dd r\wedge \alpha$ is nondegenerate on the vertical direction $\expval{\partial_\theta}$ and the $r$-direction. Hence $(1+\lambda r)\pi^*\omega_B +\dd r\wedge \alpha$ is nondegenerate whenever $1+\lambda r\neq 0$.
        \item Let $\partial_\theta\in \mathfrak{X}(Z)$ again denote the infinitesimal generator of the principal $S^1$-action on $Z$. Then
        \begin{align*}
            \iota_{\partial_\theta} \omega = (1+\lambda r)\iota_{\partial_\theta}(\pi^*\omega_B) + \iota_{\partial_\theta}(\dd r\wedge \alpha) =-\dd r,
        \end{align*}
        using that $\pi_*(\partial_\theta) = 0$, $\alpha(\partial_\theta)=1$ and $\dd r(\partial_\theta)=0$. 
        \item Since $\pi\circ i_c:r^{-1}(c)\to B$ is constant on the $S^1$-orbits, there exists a unique smooth map $\chi_c:M_c\to B$ making the diagram
    \begin{equation*}
        \begin{tikzcd}
            r^{-1}(c) \arrow[r, "i_c", hook] \arrow[d, "/S^1","\pi_c"', two heads] & {(M,\omega)} \arrow[d, "\pi", two heads] \\
            {(M_c,\omega_c)} \arrow[r, "\chi_c"]                              & {(B,\omega_B)}                               
        \end{tikzcd}
    \end{equation*}
    commute. The map $i_c$ is a diffeomorphism from $r^{-1}(c)$ to $Z \times \{c\} \subset Z \times I$, and since the $S^1$-action generated by $r$ on $r^{-1}(c)$ coincides with the principal bundle action on $Z$, we deduce that the map $\chi_c$ induced on the respective quotients is a diffeomorphism. 

    Furthermore, we compute
    \begin{align*}
        \pi_c^*\omega_c &= i_c^*((1+\lambda r)\pi^*\omega_B+\dd r\wedge\alpha) \\
        &= (1+\lambda c)(\pi\circ i_c)^*\omega_B \\
        &= (1+\lambda c)(\chi_c \circ \pi_c)^*\omega_B \\
        &= \pi_c^*\left( (1+\lambda c)\chi_c^*\omega_B\right)
    \end{align*}
    and since $\pi_c$ is a surjective submersion, the claim follows. \qedhere
    \end{enumerate}
\end{proof}


\end{document}
